% sec:schw-derivation — Derivation and Birkhoff's Theorem
\section{Derivation and Birkhoff's Theorem}
\label{sec:schw-derivation}

The overview in \cref{sec:gr-exact-solutions} stated the Schwarzschild
metric as a fait accompli; this section shows where it comes from.  The
derivation is instructive beyond its result: spherical symmetry and the
vacuum equations are so constraining that they determine the metric up
to a single parameter --- the mass.  We need not even assume that the
spacetime is static; the Einstein equations enforce staticity on their
own.

\paragraph{The spherically symmetric ansatz.}

A spacetime is spherically symmetric when the rotation group $SO(3)$
acts on it by isometries, with each orbit a two-sphere.  The most
general metric compatible with this symmetry can be written
\cite{Carroll2004}
%
\begin{equation}
  \d s^2 = -e^{2\alpha(t,r)}\,\d t^2
           + e^{2\beta(t,r)}\,\d r^2
           + r^2\!\left(\d\theta^2
             + \sin^2\!\theta\;\d\varphi^2\right) ,
  \label{eq:schw-ansatz}
\end{equation}
%
where $\alpha(t,r)$ and $\beta(t,r)$ are two unknown functions and $r$
is the \emph{areal radius}: the area of each symmetry two-sphere is
$4\pi r^2$ by construction.
\physnote{Defining $r$ so that the angular part is
$r^2\,\d\Omega^2$ is a gauge choice --- it fixes what we mean by the
radial coordinate without loss of generality.}
The exponential parametrisation keeps the metric components manifestly
nonzero and simplifies the Ricci-tensor algebra.  Note that $\alpha$
and $\beta$ are allowed to depend on both $t$ and $r$ --- we are not
assuming time-independence.  The metric therefore describes a
two-parameter family of geometries; the vacuum equations will reduce
those two functions to a single constant.

\paragraph{Solving the vacuum equations.}

The vacuum condition $\ricci = 0$ (\cref{eq:vacuum-efe}) supplies ten
equations, but spherical symmetry reduces them to four independent
components.  Computing the Christoffel symbols from
\cref{eq:schw-ansatz} and contracting --- a mechanical but lengthy
calculation --- gives the off-diagonal Ricci component
%
\begin{equation}
  R_{tr} = \frac{2\,\dot{\beta}}{r} \,,
  \label{eq:schw-ricci-tr}
\end{equation}
%
where an overdot denotes $\partial/\partial t$.  Setting
$R_{tr} = 0$ immediately yields $\dot{\beta} = 0$, so
$\beta = \beta(r)$ --- the function $\beta$ cannot depend on time.
\physnote{This is the first sign of Birkhoff's theorem: one of the
two metric functions is forced to be time-independent before we
touch the diagonal equations.}

With $\dot{\beta} = 0$, the remaining time-derivative terms ---
which enter $R_{tt}$ and $R_{rr}$ only through $\dot{\beta}$ and
$\ddot{\beta}$ --- vanish identically, leaving three diagonal
equations involving only radial derivatives (primes denote
$\partial/\partial r$).  The key structural feature is that $R_{tt}$
and $R_{rr}$ share the same second-order content, differing only
in sign and in their $1/r$ terms:
%
\begin{align}
  R_{tt} &= e^{2(\alpha-\beta)}\!\left[
    \alpha'' + {\alpha'}^{2} - \alpha'\beta'
    + \frac{2\alpha'}{r}\right] ,
  \label{eq:schw-ricci-tt} \\[4pt]
  R_{rr} &= -\alpha'' - {\alpha'}^{2} + \alpha'\beta'
    + \frac{2\beta'}{r} \,,
  \label{eq:schw-ricci-rr} \\[4pt]
  R_{\theta\theta} &= e^{-2\beta}\bigl[
    r\!\left(\beta' - \alpha'\right) - 1\bigr] + 1 \,.
  \label{eq:schw-ricci-thth}
\end{align}
%
The fourth component $R_{\varphi\varphi} =
\sin^2\!\theta\;R_{\theta\theta}$ is proportional and carries no
new information.

\paragraph{Relating $\alpha$ and $\beta$.}

The expressions for $R_{tt}$ and $R_{rr}$ look unwieldy, but a
well-chosen linear combination cancels almost everything:
%
\begin{equation}
  e^{2(\beta-\alpha)} R_{tt} + R_{rr}
  = \frac{2}{r}\!\left(\alpha' + \beta'\right) = 0 \,.
  \label{eq:schw-alpha-beta-sum}
\end{equation}
%
This forces $\alpha' = -\beta'$, and hence $\alpha(t,r) = -\beta(r) +
f(t)$ for some function $f$ of time alone.  The residual function
$f(t)$ enters only through $e^{2\alpha}\,\d t^2 = e^{-2\beta +
2f(t)}\,\d t^2$; redefining the time coordinate via $\d\bar{t} =
e^{f(t)}\,\d t$ sets $\alpha = -\beta$.  Relabelling $\bar{t} \to t$,
neither metric function depends on $t$, and no $\d t\,\d r$
cross-term appears: the metric is static.  We didn't assume
staticity --- the vacuum equations demanded it.

\paragraph{Integrating the angular equation.}

With $\alpha = -\beta$, the angular equation $R_{\theta\theta} = 0$
reduces to
%
\begin{equation}
  e^{-2\beta}\!\left(1 - 2r\beta'\right) = 1 \,.
  \label{eq:schw-ode}
\end{equation}
%
Expanding the derivative $\frac{\d}{\d r}\bigl[r\,e^{-2\beta}\bigr]
= e^{-2\beta}(1 - 2r\beta')$ shows that \cref{eq:schw-ode} is
equivalent to $\frac{\d}{\d r}\bigl[r\,e^{-2\beta}\bigr] = 1$.
Integrating once gives
%
\begin{equation}
  e^{-2\beta(r)} = 1 - \frac{C}{r} \,,
  \label{eq:schw-e2beta}
\end{equation}
%
where $C$ is a constant of integration.  Since $e^{2\alpha} =
e^{-2\beta}$, the full line element becomes
%
\begin{equation}
  \d s^2 = -\!\left(1 - \frac{C}{r}\right)\d t^2
           + \left(1 - \frac{C}{r}\right)^{\!-1}\d r^2
           + r^2\!\left(\d\theta^2
             + \sin^2\!\theta\;\d\varphi^2\right) .
  \label{eq:schw-derived}
\end{equation}
%
Substituting \cref{eq:schw-e2beta} into $R_{tt} = 0$ yields an
identity: the contracted Bianchi identity guarantees that the
remaining component is automatically satisfied once the other three
vanish.

\paragraph{Identifying the mass.}

The integration constant $C$ acquires physical meaning from the
weak-field limit.  At large $r$ the metric should reproduce Newtonian
gravity, where $g_{00} \approx -(1 + 2\Phi)$ with $\Phi = -M/r$ the
gravitational potential (\cref{eq:g00-potential}).  Matching gives
$C = 2M$, and \cref{eq:schw-derived} reduces to
\cref{eq:schwarzschild-metric}.  The entire geometry
of a spherically symmetric vacuum is encoded in a single number: the
mass~$M$.
\xrefnote{The weak-field identification of $M$ as the Newtonian mass
is developed in \cref{sec:gr-weak-field}.}

\paragraph{Birkhoff's theorem.}

The observation that the vacuum equations force staticity without it
being assumed is sufficiently important to state as a theorem.

\begin{theorem}[Birkhoff, 1923]\label{thm:birkhoff}
Every spherically symmetric solution of the vacuum Einstein equations
is static and isometric to the Schwarzschild metric
\eqref{eq:schw-derived}.
\end{theorem}

The proof is the derivation above: $R_{tr} = 0$ eliminates time
dependence in $\beta$; the combination
\eqref{eq:schw-alpha-beta-sum} ties $\alpha$ to $-\beta$ up to a
function of $t$ alone; and a time-coordinate redefinition removes that
function \cite{Carroll2004,Wald1984}.
\histnote{Birkhoff published this result in 1923
\cite{Birkhoff1923}, though Jebsen obtained the same conclusion
two years earlier.  The theorem is sometimes called the
Jebsen--Birkhoff theorem.}

The physical content is the gravitational analogue of Newton's shell
theorem.  A spherically symmetric star can pulsate, oscillate, or
undergo gravitational collapse, and provided it retains its symmetry,
the external field is unchanged --- it remains the Schwarzschild
geometry with the same $M$.
\physnote{Birkhoff's theorem forbids monopole gravitational radiation;
dipole radiation is separately forbidden by momentum conservation.
The lowest multipole that can radiate gravitationally is the
quadrupole.}
For the ray tracer, this means that photon trajectories outside a
collapsing spherical star are identical to those around a static black
hole of the same mass --- the dynamics of the interior are invisible
to external observers.

\paragraph{Scope and limitations.}

The solution has two singular surfaces: $r = 0$ and $r = 2M$.
To distinguish a genuine singularity from a coordinate artefact,
we need a quantity that all observers agree on --- a curvature
scalar.  The simplest is the Kretschmann scalar
$R_{\alpha\beta\gamma\delta}R^{\alpha\beta\gamma\delta}$, which for
the Schwarzschild metric evaluates to $48M^2/r^6$.  This diverges
at $r = 0$: the singularity there is genuine.  The singularity at $r = 2M$ is a coordinate artefact:
all curvature invariants are finite, and a freely falling observer
crosses $r = 2M$ in finite proper time without experiencing anything
locally dramatic.  Resolving this coordinate singularity --- and finding
coordinates that the ray tracer can use smoothly across the horizon ---
is the subject of the next two sections.
\warnnote{The Schwarzschild solution with $M < 0$ has a naked
singularity at $r = 0$ with no horizon.  Negative-mass solutions are
unphysical but occasionally useful as stress tests for numerical
codes.}
